\section{Introduction: when to irrigate}
Every irrigation system answers two questions: \emph{when} to apply water and \emph{how much}. The second is a matter of the root zone and the system (Chapters~2 and~7); the first is a matter of a \kw{trigger}, a rule that turns the state of the soil, the crop or the weather into a decision. This chapter studies the simplest trigger, a threshold on the suction measured by a sensor in the root zone, and asks three questions of it: how far can a soil dry before the trigger fires (the \emph{reach}), where does the water go after it is applied (the \emph{fate}), and how long does the soil remember a change of management (the \emph{memory}). The answers are laws derived in the studies on which this book is based from an ensemble of 10\,190 Daisy simulations over ten soils.

\section{Types of irrigation triggers and sensors}
Irrigation scheduling methods are grouped into four classes. \kw{Calendar} (fixed interval or fixed dose) methods ignore the state of the crop and are the oldest; they waste water when it rains and starve the crop in a hot spell. \kw{Water balance} (or budget) methods keep a running account of the depletion $D_r$ of the root zone from the weather ($ET_c$) and the water applied, and irrigate when $D_r$ reaches the readily available water; they need no sensor but drift without correction. \kw{Soil-based} methods read a sensor: a \emph{tensiometer} (a porous cup on a water-filled tube with a vacuum gauge, reading 0--80 kPa), a \emph{gypsum block} or \emph{granular matrix sensor} (electrical resistance, 0--200 kPa), or a \emph{dielectric probe} (water content). \kw{Plant-based} methods read leaf water potential, canopy temperature or stem diameter.

Soil-based triggers are of two kinds: a \emph{suction trigger} $h_t$ (irrigate when the sensor reads more than $h_t$ cm) and a \emph{depletion trigger} $f_t$ (irrigate when a fraction $f_t$ of the available water has gone). The two are linked by the retention curve of the soil, $f=f(h)$ (Equation~\ref{eq:fdep}, Chapter~9). A suction of 300 cm is a light depletion in a sand and a deep one in a silt loam; that is the reason the same trigger behaves differently in different soils.

\exbox{A suction trigger as a depletion}{For the sand S03 ($\alpha=0.145$, $n=2.68$) and the silt loam S07 ($\alpha=0.020$, $n=1.41$) find the depletion $f$ that a suction trigger of 300 cm means (field capacity at 100 cm, wilting point at 15\,849 cm).}\sol \givenline{$\alpha=0.145$ ; $n=2.68$ ; $\alpha=0.020$ ; $n=1.41$}\ans Sand: $f(300)=[S_e(100)-S_e(300)]/[S_e(100)-S_e(15849)]=(0.0112-0.0018)/(0.0112-0.0000)=0.842$; Silt loam: $f(300)=[S_e(100)-S_e(300)]/[S_e(100)-S_e(15849)]=(0.6859-0.4691)/(0.6859-0.0943)=0.366$ \ans The same 300 cm trigger is a near-complete depletion in the sand and a modest one in the silt loam. \finans

\section{The reach of a trigger}
A trigger can fire only if the soil can reach the suction at which it is set. A crop dries the soil until the conductivity around its roots falls so low that the soil can no longer supply the demand; beyond that suction the roots meet a wall. Writing that $K_r(h_c)=\kappa$ at that point, and using the dry-end form $K_r\simeq m^2(\alpha h)^{-E}$ with $E=2n+\ell(n-1)$ (Chapter~4, Equations on the reach), the reach is

\begin{equation}h_c=\frac1\alpha\Bigl(\frac{m^2}{\kappa}\Bigr)^{1/E},\qquad m=1-\frac1n,\ \kappa=3.5\times10^{-10},\ \ell=0.5.\label{eq:reach3}\end{equation}
\par\noindent where $h_c$ is the reach (cm), $\alpha$ (cm$^{-1}$) and $n$ the van Genuchten parameters, $m=1-1/n$, $\kappa$ the demand constant, $\ell$ the connectivity parameter and $E=2n+\ell(n-1)$ the exponent.



A trigger $h_t<h_c$ can fire and irrigates; a trigger $h_t>h_c$ never fires and the field behaves as rainfed. For a design, this gives an upper limit for the trigger of each soil.

\begin{table}[!htbp]\centering\small\caption{Reach of a suction trigger for the ten soils ($\kappa=3.5\times10^{-10}$, $\ell=0.5$) computed from the rounded parameters.}
\begin{tabular}{llrrrr}\toprule
\hd{Soil} & \hd{Class} & \hd{$n$} & \hd{$E$} & \hd{$h_c$ (cm)} & \hd{300 cm fires?}\\\midrule
S01 & Coarse loam & 1.35 & 2.88 & 7\,814 & yes\\
S02 & Coarse sand & 1.60 & 3.50 & 5\,359 & yes\\
S03 & Sand & 2.68 & 6.20 & 199 & no\\
S04 & Loamy sand & 2.28 & 5.20 & 425 & yes\\
S05 & Sandy loam & 1.89 & 4.22 & 1\,615 & yes\\
S06 & Loam & 1.56 & 3.40 & 9\,186 & yes\\
S07 & Silt loam & 1.41 & 3.02 & 29\,528 & yes\\
S08 & Sandy clay loam & 1.48 & 3.20 & 7\,559 & yes\\
S09 & Clay loam & 1.31 & 2.78 & 47\,608 & yes\\
S10 & Silty clay loam & 1.23 & 2.58 & 127\,830 & yes\\\bottomrule\end{tabular}\end{table}

\exbox{Reach of a loam}{Compute the reach of the loam S06 ($\alpha=0.036$, $n=1.56$) step by step, and say whether a trigger of 1000 cm would fire.}\sol \givenline{$\alpha=0.036$ ; $n=1.56$}\ans $m=1-1/1.56=0.3590$, $m^2=0.1289$; $E=2n+\ell(n-1)=3.12+0.5\times0.56=3.40$. $m^2/\kappa=0.1289/3.5\times10^{-10}=3.682e+08$; its power $1/E=0.2941$ is $330.70$, so $h_c=330.70/0.036=9186$ cm \ans A trigger of 1000 cm is below the reach and fires. \finans

\section{The fate of an irrigation millimetre: the dose law}
An irrigation millimetre does not all go to the crop. Written against a rainfed twin of the same run, the water applied $I$ is shared among transpiration, soil evaporation, other evaporation, drainage and storage; the shares $\eta_X$ (the \kw{marginal fate} in flux $X$) sum to one. The drained share depends mostly on the dose $d$ per event, because a larger dose wets deeper and exceeds the capacity of the root zone more often:

\begin{equation}\phi_D(d)=\eta_D\bigl(1-e^{-d/d_D}\bigr),\label{eq:dose3}\end{equation}


where $\eta_D$ is the saturation level of the drained share and $d_D$ the dose scale. The law is a fit to the ensemble, not a derivation; its use is to choose a dose that keeps drainage low.

\exbox{Drained share against dose}{A soil with $\eta_D=0.45$ and $d_D=9$ mm gives the dose law. Find the drained share at doses of 5, 10, 15, 20 and 30 mm per event, and the dose at which half of the saturation level is reached.}\sol \givenline{$\eta_D=0.45$ ; $d_D=9$ mm gives}\ans $d=5$: $\phi_D=0.45(1-e^{-5/9})=0.192$; $d=10$: $\phi_D=0.45(1-e^{-10/9})=0.302$; $d=15$: $\phi_D=0.45(1-e^{-15/9})=0.365$; $d=20$: $\phi_D=0.45(1-e^{-20/9})=0.401$; $d=30$: $\phi_D=0.45(1-e^{-30/9})=0.434$. Half of saturation: $1-e^{-d/9}=0.5$, $d=9\ln2=6.2$ mm \ans Doses below about 6 mm lose less than a third of the saturation level to drainage, so frequent small doses are efficient in this soil; a 30 mm dose loses almost the full 0.45. \finans

\section{Hydraulic memory}
When irrigation is changed (a mulch is added, a trigger lowered), the soil does not respond at once, because the water that was held before the change must drain or be used. This persistence is the \kw{hydraulic memory}. If a saving of $s$ millimetres is made in a flux in year~1, the \kw{memory fate} $\phi_x$ of a flux $x$ is the share of the saving that appears in $x$ in the following years; it behaves as a pooled slope of the flux change on the storage change. Large in soils that hold water (fine soils) and small in free-draining soils, it limits how much of a saving can be ``banked'' in the root zone. Chapter~7 of the application tabulates it for ten soils (Table at the end of this chapter).

